\documentclass[11pt,a4paper]{article}
\usepackage[T1]{fontenc}
\usepackage[utf8]{inputenc}
\usepackage{lmodern}
\usepackage[
margin=0.60in]{geometry}
%\linespread{1.35}
\usepackage{xcolor}
\usepackage{tabularx}
\usepackage{array}
\usepackage{enumitem}
\usepackage{titlesec}
\usepackage{microtype}
\usepackage{hyperref}
\usepackage{fancyhdr}
\usepackage{ragged2e}
\definecolor{Accent}{HTML}{1F4E79}
\definecolor{TextGray}{HTML}{444444}
\definecolor{RuleGray}{HTML}{D7DCE2}
\hypersetup{
colorlinks=true,
urlcolor=Accent,
linkcolor=Accent,
pdfauthor={Ronald Katende},
pdftitle={Ronald Katende - Academic Curriculum Vitae}
}
\pagestyle{fancy}
\fancyhf{}
\renewcommand{\headrulewidth}{0pt}
\fancyfoot[R]{\scriptsize Page \thepage}

\setlength{\parindent}{0pt}
\setlength{\parskip}{2pt}
\renewcommand{\arraystretch}{1.08}
\linespread{1.1}

\setlist[itemize]{
leftmargin=1.3em,
itemsep=1.5pt,
topsep=2pt,
parsep=0pt,
partopsep=0pt
}

\setlist[enumerate]{
leftmargin=1.55em,
itemsep=2.5pt,
topsep=2pt,
parsep=0pt,
partopsep=0pt
}

\titleformat{\section}
{\large\bfseries\color{Accent}}
{}{0pt}{}
[\vspace{-3pt}\color{RuleGray}\rule{\linewidth}{0.6pt}]

\titlespacing*{\section}{0pt}{7.5pt}{4pt}

\newcommand{\entry}[3]{%
\begin{tabularx}{\linewidth}
{@{}>{\raggedright\arraybackslash}p{0.17\linewidth} X@{}}
\textcolor{TextGray}{\textit{#1}}
&
\textit{#2}\\[-1pt]
&
#3
\end{tabularx}
\vspace{1.5pt}
}

\begin{document}

{\fontsize{21}{23}\selectfont\bfseries RONALD KATENDE, PhD}\\[-1pt]

{\large\color{Accent}\textit{
Applied Mathematics | Numerical Analysis | Scientific Computing |
Scientific Machine Learning
}}\\[4pt]

\small
Department of Mathematics, Kabale University, Uganda
\;|\;
\href{mailto:rkatende@kab.ac.ug}{rkatende@kab.ac.ug}
\;|\;
\href{mailto:rkatende92@gmail.com}{rkatende92@gmail.com}
\\[-1pt]
\href{https://ronaldkatende.github.io/}{ronaldkatende.github.io}
\;|\;
\href{https://orcid.org/0000-0002-8545-1833}
{ORCID: 0000-0002-8545-1833}
\;|\;
\href{https://scholar.google.com/citations?user=SvIwEd8AAAAJ}
{Google Scholar}

\section{Academic Profile}

Applied mathematician working in numerical analysis, scientific computing,
partial differential equations, mathematical modelling and scientific
machine learning. My research asks:
\textit{What information and computation are sufficient for a reliable scientific decision?}

I study this through stability, inverse problems, numerical linear algebra,
sparse sensing, reduced-order methods and reliable learning under limited
information and computation.

\section{Research Programme and Interests}

\textit{\textbf{Computational mathematics and numerical analysis.}}
Partial differential equations, finite-element and finite-difference methods,
numerical linear algebra, stability, conditioning, approximation and iterative
methods.

\textit{\textbf{Scientific machine learning and inverse problems.}}
Physics-informed neural networks, operator and inverse learning, sparse sensing,
low-rank and reduced-order methods, uncertainty quantification and
condition-aware reconstruction.

\textit{\textbf{Learning, computation and representation.}}
Information requirements for reliable learning and recovery; task-relative
numerical linear algebra; computational work and span; representation complexity;
and mathematical limits of learning architectures.

\textit{\textbf{Dynamical systems and scientific modelling.}}
Differential-equation models, nonlinear dynamics, epidemic systems, scientific
forecasting and computational modelling, particularly under limited data,
sensing and computational resources.

% ============================================================
\section{Academic Appointments}

\entry{2026--present}
{Lecturer, Department of Mathematics, Kabale University, Uganda}
{Undergraduate and postgraduate teaching; independent research;
postgraduate supervision; undergraduate research supervision;
curriculum development; and departmental and faculty service.}

\entry{2019--2026}
{Assistant Lecturer, Department of Mathematics, Kabale University, Uganda}
{Teaching in mathematics, statistics and scientific computing;
student research supervision; research; and postgraduate teaching
from 2025.}

\entry{2016--2019}
{Assistant Lecturer, Uganda Christian University, Uganda}
{Undergraduate teaching in mathematics, statistics, scientific
computing and quantitative methods.}

% ============================================================
\section{Education}

\entry{2025}
{PhD in Mathematics, Makerere University, Uganda}
{Specialisation in Scientific Computing and Machine Learning.
Research collaboration at the University of Bergen, Norway.
Thesis: \textit{Numerical Analysis of Physics-Informed Neural Networks
for Solving PDEs}.}

\entry{2017}
{MSc in Mathematics, Pan African University / JKUAT, Kenya}
{Dissertation:
\textit{On the Location of a Free Boundary for American Options}.}

\entry{2014}
{BSc in Mathematics and Statistics, Makerere University, Uganda}
{First Class Honours.}

\section{Peer-Reviewed and Accepted Publications}

\begin{enumerate}

\item
Katende, R.
Non-Asymptotic Stability and Consistency Guarantees for
Physics-Informed Neural Networks via Coercive Operator Analysis.
\textit{Communications in Nonlinear Science and Numerical Simulation}
159, 109909 (2026).
\href{https://doi.org/10.1016/j.cnsns.2026.109909}
{doi:10.1016/j.cnsns.2026.109909}.

\item
Katende, R. (2026).
Information contraction and stability under inverse-stable memory clocks in data-driven learning.
\textit{Research in Mathematics}, 13(1), 2743449.
\href{https://doi.org/10.1080/27684830.2026.2743449}{doi.org/10.1080/27684830.2026.2743449}

\item
Katende, R.
Rethinking Data-Efficient Artificial Intelligence for Low-Resource
Settings.
\textit{Machine Learning with Applications}
23, 100796 (2026).
\href{https://doi.org/10.1016/j.mlwa.2025.100796}
{doi:10.1016/j.mlwa.2025.100796}.

\item
Katende, R., Kasumba, H., Kakuba, G., \& Mango, J. M.
On the Error Bounds for ReLU Neural Networks.
\textit{IAENG International Journal of Applied Mathematics}
54(12), 2602--2611 (2024).
\href{https://www.iaeng.org/IJAM/issues_v54/issue_12/IJAM_54_12_10.pdf}
{Article link}.

\item
Katende, R., Kasumba, H., Kakuba, G., \& Mango, J. M.
A Proof of Convergence and Equivalence for 1D Finite Element Methods
and ReLU Neural Networks.
\textit{Annals of Mathematics and Computer Science}
25, 97--111 (2024).
\href{https://doi.org/10.56947/amcs.v25.392}
{doi:10.56947/amcs.v25.392}.

\item
Mango, J. M., Katende, R., \& Kasumba, H.
Adaptive Error-Bounded Hierarchical Matrix Compression for Efficient
and Accurate Physics-Informed Neural Networks.
\textit{Engineering Letters}
33(5), 1605--1622 (2025).
\href{https://www.engineeringletters.com/issues_v33/issue_5/EL_33_5_38.pdf}
{Article link}.

\item
Namugera, F., Katende, R., Kurama, O., \& Magero, J. M.
Long-Range Epidemic Spreading with Aging in Structured Networks.
\textit{Physica A: Statistical Mechanics and its Applications}
676, 130841 (2025).
\href{https://doi.org/10.1016/j.physa.2025.130841}
{doi:10.1016/j.physa.2025.130841}.

\item
Katende, R.
A Data-Efficient Framework for Analysing Structural Transformation
in Low- and Middle-Income Economies.
\textit{Scientific Reports}
15, 34216 (2025).
\href{https://doi.org/10.1038/s41598-025-15952-3}
{doi:10.1038/s41598-025-15952-3}.

\item
Katende, R.
Bridging the Macro--Micro Divide through a New Paradigm for Climate
Resilience Assessment in Data-Scarce Regions.
\textit{PLOS Climate}
5(1), e0000662 (2026).
\href{https://doi.org/10.1371/journal.pclm.0000662}
{doi:10.1371/journal.pclm.0000662}.

\item
Katende, R.
Modeling Resource Conflict and Inequality in Fragile Systems:
A Spatial Agent-Based Framework with a Uganda Case Study.
\textit{PLOS ONE}
21(9), e0357283 (2026).
\href{https://doi.org/10.1371/journal.pone.0357283}
{doi:10.1371/journal.pone.0357283}.

\item
Katende, R.
Cross-Country Comparative Analysis of Climate Resilience and
Localized Mapping in Data-Sparse Regions.
\textit{Frontiers in Environmental Science}
13, 1495950 (2025).
\href{https://doi.org/10.3389/fenvs.2025.1495950}
{doi:10.3389/fenvs.2025.1495950}.

\end{enumerate}

\section{Selected Conference Activity}

\begin{tabularx}{\linewidth}
{@{}>{\raggedright\arraybackslash}p{0.28\linewidth} X@{}}

\textit{DeepMath 2026}
&
Accepted poster:
\textit{When Can Depth Replace Width? A Sharp ReLU Depth--Width
Frontier with Persistent State}.
\\[3pt]

\textit{SIAM Applied Linear Algebra 2027}
&
\textit{Task-Relative Numerical Linear Algebra:
Same Spectrum, Different Work};
contributed lecture submission.
\\[3pt]

\textit{FLICS 2026}
&
Two papers accepted for presentation, Valencia, Spain.

\end{tabularx}

% ============================================================
\section{Research Funding and Funded Projects}

\begin{itemize}

\item
\textit{Kabale University Research and Innovation Fund},
Awardee, 2025--2026.

\item
\textit{Government of Uganda Research and Innovation Fund},
Co-Investigator, 2022--2024.

\item
\textit{Government of Uganda Research and Innovation Fund},
Co-Investigator, 2023--2025.


\end{itemize}

\section{Awards and Fellowships}

\begin{itemize}

\item
\textit{TWAS Research Award}, 2025.

\item
\textit{Computational Geometry Fellowship},
University of Agder, Norway, 2024.

\item
\textit{Mathematics for Africa Research Prize},
European Mathematical Society, 2023.

\item
\textit{CIMPA Fellowship}, 2023.

\end{itemize}

% ============================================================
\section{Postgraduate Research Supervision}

\begin{itemize}
\item Two MSc students successfully completed under my supervision/co-supervision;
three additional MSc researchers supervised or co-supervised in applied
mathematics, scientific computing and mathematical modelling.

\item Current postgraduate projects include differential-equation modelling,
physics-informed neural networks, inverse parameter estimation,
dynamical systems and computational mathematics.

\item Invited to contribute to PhD supervision in mathematical and
computational research projects.

\item Supervised or guided more than 50 undergraduate research projects in
mathematics, statistics, modelling and computational methods.
\end{itemize}

\section{Teaching Experience}

\begin{tabularx}{\linewidth}
{@{}>{\raggedright\arraybackslash}p{0.28\linewidth} X@{}}

\textit{Numerical \& Applied Mathematics}
&
Numerical Analysis; Advanced Numerical Methods (PG);
Differential Equations; Partial Differential Equations;
Numerical Linear Algebra; Linear Programming.
\\[3pt]

\textit{Analysis \& Linear Algebra}
&
Linear Algebra; Real Analysis; Functional Analysis.
\\[3pt]

\textit{Probability, Statistics \& Data}
&
Probability and Statistics; Statistical Modelling;
Data Analytics; Applied Data Structures.
\\[3pt]

\textit{Programming \& Scientific Computing}
&
Python Programming (UG/PG); R Programming;
scientific-computing instruction for mathematics and statistics students.

\end{tabularx}

\vspace{3pt}

\begin{itemize}

\item
Regular teaching load of more than \textit{11 credit units per semester}.

\item
Teach across undergraduate and postgraduate programmes in
mathematics, statistics, numerical methods and scientific computing.

\item
Delivered intensive Python training in Statistical Modelling
for MSc Statistics students.

\item
Delivered a one-month Python and scientific-computing course
for MSc Computational Mathematics students at Makerere University.

\item
Experience in student assessment, research-project supervision,
course development and computational laboratory instruction.

\end{itemize}

% ============================================================
\section{International Research Experience and Collaboration}

\begin{tabularx}{\linewidth}
{@{}>{\raggedright\arraybackslash}p{0.30\linewidth} X@{}}

\textit{University of Bergen, Norway}
&
Research collaboration in scientific machine learning and numerical
analysis during doctoral research.
\\[3pt]

\textit{Area Science Park, Trieste, Italy}
&
Three-month research visit in 2025 focused on sparse optimisation
and computation under data scarcity.
\\[3pt]

\textit{Justus Liebig University Giessen, Germany}
&
Invited research visit and academic seminar, 2025.
\\[3pt]

\textit{University of Havana, Cuba}
&
Invited academic visit and seminar, 2023.
\\[3pt]

\textit{University of Milan}
&
Research collaboration and academic engagement.
\\[3pt]

\textit{University of Rwanda}
&
Institutional and academic engagement involving Kabale University
and the University of Rwanda.

\end{tabularx}

% ============================================================
\section{Selected Research Prototypes and Applied Work}

\begin{tabularx}{\linewidth}
{@{}>{\raggedright\arraybackslash}p{0.23\linewidth} X@{}}

\textit{AMMED}
&
Multi-timescale memory methods for early detection of clinical
deterioration under explicit false-alarm constraints, with transfer
testing across independent hospital populations.
\\[4pt]

\textit{GridTwin Sentinel}
&
Low-compute early-warning and decision-support prototype for
electricity infrastructure combining forecasting, change detection,
temporal memory, physical consistency checks, network localisation
and uncertainty.
\\[4pt]

\textit{Industrial digital twins}
&
Sparse-sensor full-field reconstruction, condition-aware sensor
selection and cross-regime validation for high-dimensional
physical systems.
\\[4pt]

\textit{Open research practice}
&
Reproducible numerical experiments, version-controlled scientific
code and research-software development across mathematical,
machine-learning and scientific-computing projects.

\end{tabularx}

\section{Selected Research Contributions}

\textit{Scientific machine learning.}
Stability and consistency analysis for physics-informed neural networks, with links to finite-element methods.

\textit{Task-relative computation.}
Mathematical criteria for when global matrix descriptors fail to capture the work required by a specific numerical task.

\textit{Sparse recovery and sensing.}
Condition-aware reconstruction from limited measurements, including sensor selection and reduced-order methods.

\textit{Learning under limited information.}
Information, memory, computation and measurement limits for reliable learning and recovery.

\textit{Applied mathematical modelling.}
Applications to epidemics, climate, conflict, infrastructure, digital twins and clinical early warning.

% ============================================================
\section{Academic Leadership and Programme Development}

\begin{itemize}

\item
Mathematics representative on the Faculty of Science Research Committee,
Kabale University.

\item
Coordinator of Innovations and Exhibitions, Department of Mathematics, Kabale University.

\item
Developed and submitted a proposed academic programme in Computational Science and Industrial Analytics.

\item
Represented Kabale University at the National Council for Higher Education Innovation Exhibition.

\item
Contribute to curriculum development, postgraduate supervision,
departmental research development and academic service.

\end{itemize}

% ============================================================
\section{Service to the Research Community}

\textit{\textbf{Journal peer review}:}
\textit{Scientific Data};
\textit{Neurocomputing};
\textit{Machine Learning: Science and Technology};
\textit{International Journal of Machine Learning and Cybernetics};
\textit{International Journal of Data Science and Analytics};
\textit{Letters in Spatial and Resource Sciences};
\textit{Discover Artificial Intelligence};
\textit{Discover Sustainability};
\textit{Future Business Journal};
\textit{Journal of Machine Learning and Data Science}.

\vspace{3pt}

\textit{\textbf{Conference reviewing}:}
ITU Kaleidoscope;
Africa at NeurIPS 2026 Workshop
(three research submissions reviewed).

%% ============================================================
%\section{Research Metrics}
%
%Google Scholar, October 2026:
%\textit{37 citations};
%\textit{h-index 3};
%\textit{i10-index 1}.

% ============================================================
\section{Technical Competencies}

\begin{tabularx}{\linewidth}
{@{}>{\raggedright\arraybackslash}p{0.27\linewidth} X@{}}

\textit{Scientific computing}
&
Python, MATLAB, R, NumPy, SciPy, PyTorch, scikit-learn,
Git, Linux and LaTeX.
\\[3pt]

\textit{Numerical computation}
&
Numerical simulation, optimisation, numerical linear algebra,
differential-equation solvers, inverse computation and
scientific visualisation.
\\[3pt]

\textit{Research practice}
&
Reproducible computational experiments, research-software
development, version-controlled workflows and preparation of
technical and mathematical research outputs.

\end{tabularx}

% ============================================================
\section{References}

Academic referees available upon request and supplied directly
where required by the recruiting institution.

\end{document}